\frontmatter

% Voorpagina met titel
\title{Attribute-based Credentials}
\author{ir. Pim Vullers}

\maketitle

Alternatives:
\begin{itemize}
  \item Applet-based e-Ticketing
  \item Privacy for Dummies -- Smart Card Implementations of Anonymous 
Credential Systems
\end{itemize}


% Copyright informatie en ISBN (op achterzijde van de voorpagina)
%
% ISBN:
%   In het proefschrift dient in de regel een ISBN (Internationaal Standaard
%   Boeknummer) te worden aangebracht. U kunt deze aanvragen via internet op
%   http://www.isbn.nl/.
\thispagestyle{empty}

\noindent
Copyright \textcopyright\ 2014 Pim Vullers \\\null

\noindent
IPA Dissertation Series: 2014-15 \\
ISBN: 978-94-6259-376-3 \\
NUR: 980 \\\null

\noindent 
Typeset using \LaTeXe\ (\TeX Live 2013) \\
Cover design by Koen Vullers \\
Printed by Ipskamp Drukkers (\url{http://proefschriften.net/}) \\\null

\vspace{35mm}

\begin{center}
\includegraphics[scale=.7,valign=M,raise=.5ex]{images/ipalogo}
\qquad
\includegraphics[valign=M,raise=.5ex]{images/radboudlogo}
\end{center}

\noindent
The work in this thesis has been carried out under the auspices of the research school IPA (Institute for Programming research and Algorithmics).
The author was employed at the Radboud University Nijmegen.

\vspace{20mm}

\begin{center}
\includegraphics[width=.45\textwidth,valign=M,raise=.5ex]{images/tlslogo}
\qquad
\includegraphics[width=.45\textwidth,valign=M,raise=.5ex]{images/openticketinglogo}
\end{center}

\noindent 
This research was sponsored by Trans Link Systems / Open Ticketing.

\vspace{20mm}

\begin{wrapfigure}{l}{44mm}
  \vspace{-4.7mm}
  \includegraphics[scale=1.75]{images/cc-by-sa}
\end{wrapfigure}

\noindent{\small
This work is licensed under the Creative Commons Attribution-ShareAlike 4.0 
International license. 
To view a copy of this license, please visit \\
\url{http://creativecommons.org/licenses/by-sa/4.0/}.}


% Verplichte titelpagina in het Nederlands
%
% De taal van het proefschrift:
%   Het proefschrift wordt geschreven in het Nederlands, of in het Engels, het
%   Frans, het Spaans of het Duits, of - met goedkeuring van de rector
%   magnificus in een andere taal. Indien het proefschrift in het Nederlands is
%   geschreven, worden daaraan een vertaling van de titel en een samenvatting
%   van de inhoud in het Engels bijgevoegd; desgewenst kan daarnaast ook een
%   samenvatting worden bijgevoegd in een andere taal naar keuze. Indien het
%   proefschrift niet in het Nederlands is geschreven, worden in elk geval de
%   titel en een samenvatting van de inhoud in het Nederlands en ten minste in
%   ��n van de hierboven genoemde talen toevoegd. De titelpagina is altijd in
%   het Nederlands, een titelpagina in ��n van bovengenoemde talen mag worden
%   toegevoegd. (Promotiereglement, pagina 25)
%
% Titelblad:
%   Het titelblad van het proefschrift wordt als volgt ingedeeld:
%    a. de titel;
%    b. de ondertitel (eventueel);
%    c. een wetenschappelijke proeve op het gebied van ......;
%    d. de volgende tekst:
%         Proefschrift ter verkrijging van de graad van doctor aan de Radboud
%         Universiteit Nijmegen op gezag van de rector magnificus, (naam),
%         volgens besluit van het College van Decanen in het openbaar te
%         verdedigen op (datum), om (tijd) uur precies door (voornamen voluit,
%         gevolgd door de achternaam) geboren op (geboortedatum) te
%         (geboorteplaats).;
%   aan de voet van de pagina: de naam van de uitgever en de plaats van uitgave.
%   (Promotiereglement, pagina 26)
%
% Achterzijde titelblad:
%   De naam van de promotor wordt vermeld op de achterzijde van het titelblad.
%   Zijn er meerdere promotores benoemd, dan dient uit de volgorde van de namen
%   van de promotores te blijken welke hoogleraar bij de voorbereiding van het
%   proefschrift het meest betrokken is geweest. Indien een copromotor is
%   benoemd dan wordt diens naam ook op de achterzijde van het titelblad vermeld
%   en wanneer er meerdere copromotores zijn benoemd dan dient eveneens uit de
%   volgorde van hun namen te blijken wie het meest bij de voorbereiding van het
%   proefschrift betrokken is geweest. Indien promotor(es) of copromotor(es) van
%   elders komen, dient de universiteit van herkomst te worden vermeld.
%   Op deze pagina worden tevens de namen vermeld van de leden van de
%   manuscriptcommissie. Indien een lid van de manuscriptcommissie van elders
%   komt, dient de universiteit of het instituut van herkomst te worden vermeld.
%   (Promotiereglement, pagina 26)
\thispagestyle{empty}

\begin{center}
  \textbf{\Large Applet-based e-Ticketing}\\[15mm]

  Een wetenschappelijke proeve op het gebied van de \\
  Natuurwetenschappen, Wiskunde en Informatica \\[15mm]

  \textsc{Proefschrift} \\[15mm]

  ter verkrijging van de graad van doctor \\
  aan de Radboud Universiteit Nijmegen \\
  op gezag van de rector magnificus, prof. mr. S.C.J.J. Kortmann, \\
  volgens besluit van het College van Decanen \\
  in het openbaar te verdedigen op \textbf{vrijdag 16 augustus 2013} \\
  om \textbf{..:.. uur} precies. \\[30mm]

  door \\[30mm]

  Pim Vullers \\[15mm]

  geboren op 20 juni 1986 \\
  te Venlo
\end{center}

\clearpage

\begin{flushleft}
  \thispagestyle{empty}

  \textbf{Promotor:} \\
  \indent Prof. dr. B.P.F. Jacobs \\[15mm]

  \textbf{Manuscriptcommissie:}
\end{flushleft}

% Verplichte samenvating in het Nederlands:
%
% Zie 'De taal van het proefschrift' hierboven
\chapter*{Samenvatting}
\addcontentsline{toc}{chapter}{Samenvatting}

% Voorwoord in het Nederlands (hier vermelden dat de scriptie verder Engels is)
%
% Dankbetuigingen:
%   Personen en instellingen die ide�le of materi�le medewerking hebben verleend
%   aan de totstandkoming van het proefschrift of een onderdeel daarvan, kunnen
%   worden vermeld:
%    - in een voorbericht of op de achterzijde van de titelpagina indien de
%      medewerking op het gehele proefschrift betrekking heeft gehad;
%    - dan wel in een voetnoot op de eerste bladzijde van het desbetreffende
%      onderdeel indien het medewerking aan een specifiek onderdeel betreft.
\chapter*{Voorwoord}
\addcontentsline{toc}{chapter}{Voorwoord}

Het is inmiddels al weer vijf jaar geleden dat ik in Nijmegen ben gestart met 
mijn promotie traject. Aan de ene kant is het een traject dat je zelf moet 
doorlopen, maar aan de andere kant is het ook zeker een traject dat je niet 
alleen kunt doorlopen. Daarom wil ik bij deze iedereen bedanken die me tijdens 
deze periode geholpen heeft. Ik wil niemand te kort doen, maar toch wil ik een
aantal personen in het bijzonder noemen.

Om te beginnen mijn promotor, Bart, die tevens de taak van dagelijks begeleider
op zich genomen heeft. Niet alleen heb je me ondersteund bij het uitvoeren van 
mijn onderzoek en het schrijven van mijn artikelen, maar ook met het promoten 
van mijn resultaten en het opdoen van ervaring binnen het onderwijs. Ondanks je
drukke agenda was er altijd wel tijd om mijn dagelijkse beslommeringen te 
bespreken. Bart, bedankt.

Ook de artikelen waarop dit proefschrift gebaseerd is heb ik niet alleen 
geschreven. Ik wil dan ook mijn collega's, Wojciech, Gergely, Bart, Jaap-Henk
en Lejla van harte bedanken voor hun hulp. Wojciech en Erik wil ik ook graag 
bedanken voor hun kennis en ervaring met het programmeren van smartcards waar
ik veel van geleerd heb, en in dit werk veelvuldig toegepast heb.

Daarnaast wil ik ook mijn mede-promovendi bedanken voor de leuke tijd die ik 
gehad heb. Bij de promovendi school IPA was het altijd wel gezellig, in het 
bijzonder het gezelschap van Carst en Jeroen, mijn medestudenten van de TU/e,
en Wouter en Gergely, mijn collega's uit Nijmegen, en natuurlijk de organisatie
van IPA, Meivan en Tim. Ook heb ik mijn volledige promotie op een ruime kamer
gezwoegd met vele andere promovendi. Uit de gezamenlijke arbeid heb ik veel 
energie geput, maar zeker ook uit onze gesprekken en werkonderbrekingen. In het
bijzonder staan de kleine security studie projectjes samen met Gerhard en Joeri
me nog goed bij.

Tijdens mijn promotie heb ik ook de mogelijkheid gehad om bij Microsoft stage 
te lopen. Zo heb ik ervaring op kunnen doen en in de keuken kunnen kijken bij 
een onderzoeksafdeling bij een groot bedrijf. Christian, bedankt voor het
regelwerk om de stage mogelijk te maken en je begeleiding tijdens mijn tijd bij
Microsoft.

Uit mijn onderzoek is het IRMA project voortgekomen dat verder onderzoek doet
naar het gebruik van attributen in de praktijk. Er is veel werk in gaan zitten,
en ik ben blij dat ik dit niet alleen heb hoeven doen. Wouter, Gergely, Ronny,
Roland, Maarten, Jaap-Henk, Antonio, Bart en Martijn, bedankt.

Ook wil ik mijn commissieleden bedanken voor de tijd die ze gestoken hebben in 
het beoordelen van mijn proefschrift. In het bijzonder wil ik Sjouke bedanken, 
omdat je ook al mijn begeleider bent geweest tijdens mijn afstudeerproject in 
Luxemburg.

Natuurlijk heeft al dit werk ook zijn invloed gehad op mijn persoonlijke leven.
Tijdens mijn promotie zullen mijn broers Bas en Koen als mijn paranimfen 
optreden. Koen heeft ook de tijd genomen om een omslag te ontwerpen voor dit
proefschrift, en me daar veel werk mee uit handen genomen. Van mijn ouders,
Jan en Gerrie, heb ik altijd de steun gehad voor deze promotie, ook al was het
niet altijd te begrijpen waar ik dan precies mee bezig was. 

Als laatste wil ik Marlou bedanken, die het niet altijd even makkelijk heeft 
gehad tijdens mijn promotie (in het bijzonder tijdens het schrijven van dit 
proefschrift), maar me toch altijd ondersteund heeft om dit werk af te ronden. 
Bedankt.

\begin{flushright}
  \textit{Pim Vullers\\ Nooitgedacht, oktober 2014}
\end{flushright}

\cleardoublepage


% Titelpagina in het Engels
%
% Zie 'De taal van het proefschrift' hierboven.
\thispagestyle{empty}

\begin{center}
  ~ \\[20mm]

  \textbf{\Large Efficient Implementations of\\ Attribute-based Credentials on Smart Cards}\\[15mm]

  ~ \\
  An academic essay in Science \\[15mm]

  \textsc{Doctoral Thesis} \\[15mm]

  to obtain the degree of doctor \\
  from the Radboud University Nijmegen \\
  on the authority of the Rector Magnificus, prof. mr. S.C.J.J. Kortmann, \\
  according to the decision of the Council of Deans \\
  to be defended in public on \textbf{Friday, November 28, 2014} \\
  at precisely \textbf{10:30 hours}. \\[20mm]

  by \\[20mm]

  ir. Pim Vullers \\[15mm]

  born on June 20, 1986 \\
  in Venlo, The Netherlands
\end{center}

\clearpage

\thispagestyle{empty}

\noindent\textbf{Supervisor:} \\[2mm]
\indent\begin{tabular}{p{40mm}l}
  Prof. dr. B.P.F. Jacobs & Radboud University Nijmegen \\
\end{tabular} \\[2mm]

\noindent\textbf{Doctoral Thesis Committee:} \\[2mm]
\indent\begin{tabular}{p{40mm}l}
  Prof. dr. E.R. Verheul     & Radboud University Nijmegen ~--~ The Netherlands \\
  Prof. dr. S. Mauw          & University of Luxembourg ~--~ Luxembourg \\
  Prof. dr. ir. B. de Decker & KU Leuven ~--~ Belgium \\
  Dr. G. Neven               & IBM Research Z\"{u}rich ~--~ Switzerland \\
  Dr. ir. E. Poll            & Radboud University Nijmegen ~--~ The Netherlands \\
\end{tabular}


% Samenvatting in het Engels
\chapter*{Abstract}
\addcontentsline{toc}{chapter}{Abstract}

[TODO: Abstract in English.]

% Voorwoord in het Engels
%
% Zie 'Dankbetuigingen' hierboven.
\chapter*{Preface}
\addcontentsline{toc}{chapter}{Preface}

Five years have passed since I started my PhD programme in Nijmegen. On the one
hand this is a journey that one has to perform alone, on the other hand it is a
journey one cannot finish alone. Therefore I would like to say thank you to 
everybody who helped me during this period. I don't want to deprive anybody, 
but there are some that I want to thank explicitly.

First, my supervisor, Bart, who supervised my PhD thesis, but also my daily
activities. You have supported me in my research and the writing of my articles,
but you have also promoted my results and allowed me te gain experience in 
teaching. Despite your busy schedule, there has always been time to discuss my 
daily worries. Bart, thanks.

I did not write the articles, on which this PhD thesis is based, alone. Hence,
I would like to thank my co-authors, Wojciech, Wojciech, Gergely, Bart, 
Jaap-Henk and Lejla sincerely for their help. I also want to thank Wojciech and
Erik for teaching me their knowledge and experience in programming smart cards,
which has been useful throughout my research.

Daarnaast wil ik ook mijn mede-promovendi bedanken voor de leuke tijd die ik 
gehad heb. Bij de promovendi school IPA was het altijd wel gezellig, in het 
bijzonder het gezelschap van Carst en Jeroen, mijn medestudenten van de TU/e,
en Wouter en Gergely, mijn collega's uit Nijmegen, en natuurlijk de organisatie
van IPA, Meivan en Tim. Ook heb ik mijn volledige promotie op een ruime kamer
gezwoegd met vele andere promovendi. Uit de gezamenlijke arbeid heb ik veel 
energie geput, maar zeker ook uit onze gesprekken en werkonderbrekingen. In het
bijzonder staan de kleine security studie projectjes samen met Gerhard en Joeri
me nog goed bij.

Tijdens mijn promotie heb ik ook de mogelijkheid gehad om bij Microsoft stage 
te lopen. Zo heb ik ervaring op kunnen doen en in de keuken kunnen kijken bij 
een onderzoeksafdeling bij een groot bedrijf. Christian, bedankt voor het
regelwerk om de stage mogelijk te maken en je begeleiding tijdens mijn tijd bij
Microsoft.

Uit mijn onderzoek is het IRMA project voortgekomen dat verder onderzoek doet
naar het gebruik van attributen in de praktijk. Er is veel werk in gaan zitten,
en ik ben blij dat ik dit niet alleen heb hoeven doen. Wouter, Gergely, Ronny,
Roland, Maarten, Jaap-Henk, Antonio, Bart en Martijn, bedankt.

Ook wil ik mijn commissieleden bedanken voor de tijd die ze gestoken hebben in 
het beoordelen van mijn proefschrift. In het bijzonder wil ik Sjouke bedanken, 
omdat je ook al mijn begeleider bent geweest tijdens mijn afstudeerproject in 
Luxemburg.

Natuurlijk heeft al dit werk ook zijn invloed gehad op mijn persoonlijke leven.
Tijdens mijn promotie zullen mijn broers Bas en Koen als mijn paranimfen 
optreden. Koen heeft ook de tijd genomen om een omslag te ontwerpen voor dit
proefschrift, en me daar veel werk mee uit handen genomen. Van mijn ouders,
Jan en Gerrie, heb ik altijd de steun gehad voor deze promotie, ook al was het
niet altijd te begrijpen waar ik dan precies mee bezig was. 

Als laatste wil ik Marlou bedanken, die het niet altijd even makkelijk heeft 
gehad tijdens mijn promotie (in het bijzonder tijdens het schrijven van dit 
proefschrift), maar me toch altijd ondersteund heeft om dit werk af te ronden. 
Bedankt.

\begin{flushright}
  \textit{Pim Vullers\\ Nooitgedacht, October 2014}
\end{flushright}


% Inhoudsopgave
\tableofcontents